\subsection{CT5 --- Local-to-global bookkeeping}\label{ct:5}\label{the-local-versus-global-rationale}
\ctsignature{\(P_{1\to5}\to \ctcert{C4};\ \ctint{P_{5\to4},P_{5\to11},P_{5\to14}};\ \ctrec{\ctnone}\).}

CT5 is an exact sequential machine for converting a local witness family into
a global comparison.  Its \texttt{Framework} contains the complete CT4
framework together with dependent site and local-witness types, the witness
extractor predicates, locality, ledger readiness, summation, and the CT11 and
CT14 datum types.  A validated \texttt{Input} contains
\(\langle G,\mathfrak B,w\rangle\), where \(w\) is the explicit dependent
witness extractor.  The CT4 framework is embedded so that the CT5--CT4
handoff has a definitionally fixed consumer vocabulary.

The core plan has five independent fields:
\texttt{scope}, \texttt{locality}, \texttt{deficit},
\texttt{summation}, and \texttt{comparison}.  Scope admission returns either a
typed obstruction or \texttt{ScopedState}.  The locality node is a
single-successor certification boundary: \texttt{LocalityCertificate} proves
that the global estimates are supported by the declared local extractor.
The deficit node then returns a CT11 additive-deficit payload or
\texttt{LocalLedgerState}.  Only a local ledger reaches the summation node,
which certifies \texttt{SummationState} before the global comparison is
classified.

The comparison node returns exactly one of a C4 certificate, a CT4 charge
ledger, or a CT14 aggregate residual.  \texttt{CT4Payload} contains only the
assignment, capacity, and lower-bound fields that CT4 still needs to certify.
Its \texttt{toInput} adapter fixes \texttt{G} and \texttt{branch}
definitionally.  The theorems \texttt{ct4\_payload\_same\_ambient} and
\texttt{ct4\_payload\_same\_branch} therefore hold by reflexivity.  This
prevents the handoff from changing the object or branch state during ledger
assembly.

Every terminal result retains the exact locality, ledger, and summation
indices that justify it.  \texttt{Port} and \texttt{HandoffPlan} certify only
consumer acceptance.  Theorems \texttt{run\_verified}, \texttt{run\_total},
and \texttt{run\_trace\_valid} establish terminal soundness, relative
totality, and trace validity.

The soundness contracts are assigned as follows.
\begin{itemize}
\tightlist
\item \textbf{S-Def} fixes the sites, extractor, witness relation, and downstream vocabulary.
\item \textbf{S-Equiv} is the locality certification that connects each global estimate to its local witness data.
\item \textbf{S-Dich} is discharged by the scope, deficit, and comparison constructors.
\item \textbf{S-Comp} is discharged by \texttt{SummationState} and the final comparison result.
\item \textbf{S-Pers} is enforced by the dependent indices carried through the ledger, summation, and terminals.
\item \textbf{S-Rout} and \textbf{S-Trig} are discharged by the CT4, CT11, and CT14 payload constructors and the separate handoff plan.
\end{itemize}

\begin{figure}[p]
\centering
\vspace*{-1.70cm}
\resizebox{0.96\textwidth}{!}{%
\begin{tikzpicture}[x=1cm,y=1cm,every node/.style={font=\scriptsize}]
\node[bcwide] (entry) at (0,0) {{\tiny\ttfamily CT5.entry}\\\textbf{Validated local family}\\\((G,\mathfrak B,w)\)};
\node[bcdec,bclemma,text width=3.45cm,minimum width=4.15cm] (scope) at (0,-2.35) {{\tiny\ttfamily CT5.decide.scope}\\\textbf{Bounded witnesses available?}};
\node[bcscopefail,text width=3.25cm,minimum width=3.25cm] (scopeout) at (7.35,-2.35) {{\tiny\ttfamily CT5.terminal.scope}\\\textbf{Scope exit}};
\node[bcwide,bclemma] (locality) at (0,-4.80) {{\tiny\ttfamily CT5.certify.locality}\\\textbf{Locality certificate}\\each estimate has declared witness support};
\node[bcroutechoice,bctheorem,text width=3.25cm,minimum width=4.10cm] (deficit) at (0,-7.35) {{\tiny\ttfamily CT5.decide.deficit}\\\textbf{Local ledger or additive deficit?}};
\node[bcroute,bcrouteneutral,bcoutputroute,text width=3.25cm,minimum width=3.25cm] (ct11) at (-7.00,-7.35) {{\tiny\ttfamily CT5.terminal.ct11}\\\textbf{CT 11 exit}\\additive deficit};
\node[bcwide,bclemma] (summation) at (0,-10.05) {{\tiny\ttfamily CT5.certify.summation}\\\textbf{S-Comp certificate}\\local inequalities sum globally};
\node[bcroutechoice,bctheorem,text width=3.25cm,minimum width=4.10cm] (comparison) at (0,-12.75) {{\tiny\ttfamily CT5.decide.comparison}\\\textbf{Classify global comparison}};
\node[bcclosed,text width=2.90cm,minimum width=2.90cm] (c4) at (-6.60,-15.85) {{\tiny\ttfamily CT5.terminal.c4}\\\textbf{C4}};
\node[bcroute,bcrouteneutral,bcoutputroute,text width=3.25cm,minimum width=3.25cm] (ct4) at (0,-15.85) {{\tiny\ttfamily CT5.terminal.ct4}\\\textbf{CT 4 exit}\\exact ledger};
\node[bcroute,bcrouteneutral,bcoutputroute,text width=3.25cm,minimum width=3.25cm] (ct14) at (6.60,-15.85) {{\tiny\ttfamily CT5.terminal.ct14}\\\textbf{CT 14 exit}\\aggregate residual};
\draw[bcarr] (entry)--(scope);
\draw[bcscopearr] (scope)--node[bclabel,above]{exit}(scopeout);
\draw[bcarr] (scope)--node[bclabel,right]{ready}(locality);
\draw[bcarr] (locality)--node[bclabel,right]{certified}(deficit);
\draw[bchandoff] (deficit)--node[bclabel,above]{\(P_{5\to11}\)}(ct11);
\draw[bcarr] (deficit)--node[bclabel,right]{local ledger}(summation);
\draw[bcarr] (summation)--node[bclabel,right]{summed}(comparison);
\draw[bcarr] (comparison)--node[bclabel,above,sloped]{capacity gap}(c4);
\draw[bchandoff] (comparison)--node[bclabel,right]{\(P_{5\to4}\)}(ct4);
\draw[bchandoff] (comparison)--node[bclabel,above,sloped]{\(P_{5\to14}\)}(ct14);
\end{tikzpicture}%
}
\caption[Closure-tactic 5: exact local-to-global machine]{Closure-tactic 5 as an exact sequential machine.  Locality and summation are one-successor certification boundaries.  The comparison node can close at C4 or emit one indexed payload.  The CT4 payload converts to the exact downstream \texttt{CT4.Input} while preserving the ambient object and branch state definitionally.  Monospaced identifiers are shared by Lean, JSON, and executable traces.}
\label{fig:ct5-local-global-bookkeeping}
\end{figure}
\FloatBarrier
